\section{Introduction}
Particle decay tests a persistence model beyond assigning a stationary mass. A model must specify allowed products, conserved quantities, transition rates and the statistical law connecting those rates to observations. We investigate decay as a change of routing protocol in the proposed $E_8$ substrate.

The identification is made in one direction only. We first establish that the $E_8$ substrate \emph{contains} an object of a given type---a chiral grading, a threefold colour structure, a charge spectrum, four routing protocols, a tier ladder---and only then observe that the Standard-Model lexicon has a counterpart. The counterpart is a proposed correspondence that helps readers who already know the physics; it is not used as evidence that the substrate object exists. The existence questions are settled on the substrate side, and where a step is not yet proved we say so rather than importing the physical answer.

The geometric ledger proposes masses and channel weights; a transition relation specifies allowed routes; and a stochastic model assigns waiting times. These are distinct inputs. Throughout, $\tau$ denotes a mean lifetime, with $t_{1/2}=\tau\log 2$ for exponential survival.
